\makeatletter\def\input@path{{../template/vsp-beamer}{../template/minimus}}\makeatother
\documentclass[red,tutorial,sectionpages,shtu]{vsp-beamer}
\graphicspath{{../template/vsp-beamer/assets/}}
\usepackage{minimus-section}
\usepackage{minimus-text}
\usepackage{minimus-math}
\usepackage{minimus-float}
\usepackage{minimus-code}
\usepackage{minimus-colorbox}
\usepackage{minimus-reference}
\usepackage{tikz}
\usetikzlibrary{arrows.meta,positioning,calc}
\lstset{basicstyle=\ttfamily\footnotesize,numberstyle=\tiny\color{gray},breaklines=true,xleftmargin=8pt}
\title{CS100 Recitation 7}
\subtitle{C++ Class基础：封装、接口与对象生命周期}
\VSPsetspeaker{Xiangkun Meng}{CS100 Teaching Assistant}
\VSPsetupfootline{CS100 Recitation 7}{Fall 2026}{Xiangkun Meng}
\begin{document}
    \VSPtitleframe
    \begin{frame}{今天用一个Student回答三个问题}
        \begin{enumerate}
            \item \textbf{对象与接口：}学生的信息属于谁？谁可以修改？
            \item \textbf{查看与const：}只查看信息，应该作出什么承诺？
            \item \textbf{初始化与生命周期：}对象如何开始，又在何时结束？
        \end{enumerate}
        \medskip
        \begin{vspquotered}{贯穿案例}
            姓名\texttt{m\_name}、学号\texttt{m\_id}、入学年份\texttt{m\_entranceYear}。\\
            我们希望姓名和学号非空，入学年份在1900至2100之间。
        \end{vspquotered}
        \medskip
        全部代码使用C++17；片段按页逐步补齐。
    \end{frame}
    \section{对象与接口}

\begin{frame}{Class是类型，Object是具体的一个}
    \begin{columns}[T]
        \begin{column}{0.54\textwidth}
            \begin{Code}*{cpp}
                \lstinputlisting{code/members.inc}
            \end{Code}
        \end{column}
        \begin{column}{0.42\textwidth}
            \textbf{Class：}规定学生有何信息、能做何操作。\\[8pt]
            \textbf{Object：}按这个类型创建的某个学生，如Alice、Bob。\\[8pt]
            \textbf{成员变量：}每个对象各自保存的信息。\\[8pt]
            \textbf{成员函数：}对这个对象执行的操作。
        \end{column}
    \end{columns}
    \medskip
    \small 使用\texttt{std::string}需包含\texttt{<string>}。类定义最后的分号不能省。
\end{frame}

\begin{frame}{看到一个名字，先分清它从哪里来}
    \begin{Code}*{cpp}
        \lstinputlisting{code/names.inc}
    \end{Code}
    \begin{itemize}
        \item \texttt{m\_entranceYear}：\textbf{成员变量}，保存在当前学生对象里。
        \item \texttt{year}：\textbf{参数}，由这一次调用传入。
        \item \texttt{years}：\textbf{局部变量}，进入函数后计算，函数结束即离开作用域。
    \end{itemize}
    \medskip
    \alert{成员变量用\texttt{m\_}前缀，参数和局部变量不用；成员函数可直接使用自己的成员。}\\
    \small 成员函数用小驼峰命名：首词小写，后续词首字母大写，如\texttt{setName}、\texttt{getName}。
\end{frame}

\begin{frame}{public、private、protected分别给谁用？}
    \begin{columns}[T]
        \begin{column}{0.47\textwidth}
            \begin{itemize}
                \item \texttt{public}：普通类外代码也可访问。
                \item \texttt{private}：类的成员可访问，派生类不能直接访问。
                \item \texttt{protected}：类的成员、派生类可按规则访问，普通类外代码不能访问。
            \end{itemize}
            \begin{Code}*{cpp}
                \lstinputlisting{code/access.inc}
            \end{Code}
            \small 左侧仍是Student；右侧用最小继承例子单独说明protected。
        \end{column}
        \begin{column}{0.49\textwidth}
            \begin{Code}*{cpp}
                \lstinputlisting{code/protected.inc}
            \end{Code}
        \end{column}
    \end{columns}
    \medskip
    \small 权限标签持续生效到下一个标签或类结束；友元另行授权，本课不展开。\\
    \alert{派生类的函数能读取继承的protected成员，普通外部代码不能直接读写它。}
\end{frame}

\begin{frame}{课堂提问：struct与class究竟差在哪里？}
    \renewcommand{\arraystretch}{1.35}
    \begin{tabular}{lll}
        & \textbf{struct} & \textbf{class}\\
        \hline
        未写权限标签的成员 & \texttt{public} & \texttt{private}\\
        未写继承权限时 & \texttt{public}继承 & \texttt{private}继承
    \end{tabular}
    \medskip
    \begin{Code}*{cpp}
        \lstinputlisting{code/default-access.inc}
    \end{Code}
    \small \texttt{Base}已在上方定义；冒号表示继承，这里只辨认未写权限时的默认值。\\
    \alert{两者都能有成员函数、构造函数和访问权限；不是“一个只能放数据”。}
\end{frame}

\begin{frame}{public成员变量：检查可以被绕过}
    \begin{columns}[T]
        \begin{column}{0.52\textwidth}
            \begin{Code}*{cpp}
                \lstinputlisting{code/public-data.inc}
            \end{Code}
        \end{column}
        \begin{column}{0.44\textwidth}
            这段\textbf{反例}可以编译，却出现不合理年份。\\[8pt]
            即使另写一个带检查的函数，外部仍能直接赋值。\\[8pt]
            \textbf{封装：}隐藏内部表示，让对外操作维护对象的约束。
        \end{column}
    \end{columns}
    \medskip
    \begin{vspquotered}{Student的设计选择}
        把成员变量设为private，只公开需要的操作。\textbf{private是起点，接口检查才负责守住规则。}
    \end{vspquotered}
\end{frame}

\begin{frame}{接口先检查，再修改对象状态}
    \begin{columns}[T]
        \begin{column}{0.56\textwidth}
            \begin{Code}*{cpp}
                \lstinputlisting{code/set-name.inc}
            \end{Code}
        \end{column}
        \begin{column}{0.40\textwidth}
            \textbf{规则：姓名不能为空。}\\[8pt]
            输入为空：返回\texttt{false}，原姓名保留。\\[8pt]
            输入有效：修改\texttt{m\_name}，返回\texttt{true}。\\[8pt]
            \texttt{const std::string \&name}：只读取传入的字符串。
        \end{column}
    \end{columns}
    \medskip
    \alert{接口不是把赋值换个位置；它规定允许做什么，以及失败时保留什么。}
\end{frame}

\begin{frame}{跟着执行：改名成功与改名失败}
    \begin{Code}*{cpp}
        \lstinputlisting{code/use-interface.inc}
    \end{Code}
    \begin{enumerate}
        \item 创建学生，开始时姓名为\texttt{Alice}；构造过程稍后讲。
        \item \texttt{ok1 == true}：姓名变为\texttt{Alicia}。
        \item \texttt{ok2 == false}：拒绝空姓名，仍然是\texttt{Alicia}。
        \item 通过查看接口\texttt{getName()}取得姓名，最终输出\texttt{Alicia}。
    \end{enumerate}
    \alert{提问：若\texttt{m\_name}仍然public，空姓名检查还能保证规则吗？}
\end{frame}

\begin{frame}{硬件类比：我们究竟开放了哪些操作？}
    \begin{columns}[T]
        \begin{column}{0.47\textwidth}
            \begin{vspquotered}{单片机开发板}
                开放GPIO、芯片引脚等低层连接。\\[7pt]
                开发者自由度高，也需要自己遵守电压、接线等限制。\\[7pt]
                直接接触内部，影响内部状态的路径更多。
            \end{vspquotered}
        \end{column}
        \begin{column}{0.47\textwidth}
            \begin{vspquotered}{日常使用的笔记本}
                通常通过USB、键盘、电源键操作，通过屏幕查看结果。\\[7pt]
                使用者通过约定接口完成任务，无需直接操作内部芯片。\\[7pt]
                类似通过成员函数使用Student。
            \end{vspquotered}
        \end{column}
    \end{columns}
    \medskip
    \small 开发板开放引脚是有意的设计，不代表质量差。这里类比的是\textbf{接口边界与使用责任}。\\
    \alert{回到Student：隐藏成员变量，开放经过检查的操作。}
\end{frame}

    \section{查看与const}

\begin{frame}{this指向本次调用的那个对象}
    \begin{Code}*{cpp}
        \lstinputlisting{code/this.inc}
    \end{Code}
    \begin{itemize}
        \item 假定\texttt{a}在2026年入学，\texttt{b}在2028年入学：上面两次调用分别得到\texttt{true}和\texttt{false}。
        \item 在普通非静态成员函数中，\texttt{m\_entranceYear}相当于\texttt{this->m\_entranceYear}。
        \item 调用传入年份，也隐含指定了\textbf{哪个对象}；不用自己把\texttt{this}写进参数表。
    \end{itemize}
    \alert{同一个成员函数可以服务多个对象，每次读取当前对象自己的数据。}
\end{frame}

\begin{frame}{只查看对象的函数，要在末尾写const}
    \begin{Code}*{cpp}
        \lstinputlisting{code/const-view.inc}
    \end{Code}
    \begin{itemize}
        \item 参数表后的\texttt{const}约束\textbf{当前对象}：本例不能修改成员变量。
        \item 这时\texttt{this}的类型为\texttt{const Student *}，不能通过它调用非const成员函数。
        \item \texttt{getName()}返回姓名的一个值；外部修改这个返回值不会修改学生内部的姓名。
    \end{itemize}
    \alert{查看功能应声明为const，让普通对象和const对象都能使用。}
\end{frame}

\begin{frame}{少写const，函数体不改数据也会编译失败}
    \begin{Code}*{cpp}
        \lstinputlisting{code/const-error.inc}
    \end{Code}
    \begin{vspquotered}{原因与修正}
        编译器依据函数声明判断能否调用，不因函数体“看起来只读”而放行。\\[5pt]
        修正：\texttt{std::string getName() const \{ return m\_name; \}}
    \end{vspquotered}
    \small 参数中的\texttt{const}保护参数；函数末尾的\texttt{const}保护当前对象，位置与含义不同。
\end{frame}

\begin{frame}{一起判断：哪些函数应为const？}
    \begin{Code}*{cpp}
        \lstinputlisting{code/const-exercise.inc}
    \end{Code}
    \renewcommand{\arraystretch}{1.2}
    \begin{tabular}{lll}
        \textbf{调用者} & \textbf{const成员函数} & \textbf{非const成员函数}\\
        \hline
        普通\texttt{Student}对象 & 可以 & 可以\\
        \texttt{const Student}对象 & 可以 & 不可以
    \end{tabular}
    \par\medskip
    \alert{A可以调用，结果为true；B若取消注释会编译失败。修改姓名不属于查看。}
\end{frame}

    \section{初始化与生命周期}

\begin{frame}{构造函数：让对象从有效状态开始}
    \begin{Code}*{cpp}
        \lstinputlisting{code/constructor.inc}
    \end{Code}
    \begin{itemize}
        \item 名字与类相同，\textbf{没有返回类型}；创建对象时自动调用。
        \item 冒号后的\textbf{初始化列表}直接初始化成员，随后才执行函数体。
        \item 无效输入抛出异常：本次构造失败，不会得到一个构造完成的无效Student。
    \end{itemize}
    \small 需包含\texttt{<stdexcept>}；这里只把\texttt{throw}读作“报告错误并中止本次构造”，不展开异常处理。
\end{frame}

\begin{frame}{先初始化成员，再执行构造函数体}
    \begin{columns}[T]
        \begin{column}{0.56\textwidth}
            \begin{Code}*{cpp}
                \lstinputlisting{code/init-order.inc}
            \end{Code}
        \end{column}
        \begin{column}{0.40\textwidth}
            \textbf{实际顺序看声明位置：}\\[6pt]
            1. \texttt{m\_name}\\
            2. \texttt{m\_id}\\
            3. \texttt{m\_entranceYear}\\
            4. 构造函数体\\[8pt]
            调换初始化列表的顺序，\textbf{不会改变}这个顺序。
        \end{column}
    \end{columns}
    \medskip
    若只在函数体写\texttt{m\_name = name;}，字符串先默认初始化，再被赋值。\\
    未指定初值的内置整数成员可能没有确定值，不能先读后写。\\
    \alert{初始化列表与声明顺序保持一致；不要依赖尚未初始化的后续成员。}
\end{frame}

\begin{frame}{默认构造：无参数也要有合理含义}
    \begin{Code}*{cpp}
        \lstinputlisting{code/default-ctor.inc}
    \end{Code}
    \begin{itemize}
        \item 默认构造函数可以无实参调用；自己声明构造函数后，编译器不再自动提供无参版本。
        \item 本课Student要求姓名、学号和入学年份有效，因此\textbf{创建时必须提供这三项信息}。
        \item 无实参调用是否合理，取决于类的设计；不要为了让代码通过而随意补一个默认学生。
    \end{itemize}
\end{frame}

\begin{frame}{析构函数：对象结束时自动执行}
    \begin{Code}*{cpp}
        \lstinputlisting{code/lifetime.inc}
    \end{Code}
    \begin{itemize}
        \item 析构函数名为\texttt{\textasciitilde Student}，没有返回类型，也没有参数。
        \item 这里两个局部对象离开花括号时自动析构，\textbf{后创建的先结束}。
        \item 输出依次为\texttt{Bob leaves}、\texttt{Alice leaves}；即使无显式析构函数，成员仍会自动清理。
    \end{itemize}
    \alert{这里打印仅为了观察生命周期；本课Student不需要手写清理字符串。}
\end{frame}

\begin{frame}{new创建对象，delete结束动态对象}
    \begin{Code}*{cpp}
        \lstinputlisting{code/new-delete.inc}
    \end{Code}
    \begin{enumerate}
        \item \texttt{new Student(...)}分配存储，并调用构造函数；\texttt{p}得到对象地址。
        \item \texttt{p->getName()}通过指针访问成员，等价于\texttt{(*p).getName()}。
        \item \texttt{delete p}调用析构函数，然后释放存储；不能再访问原对象。
        \item \texttt{p = nullptr}只清空这个指针，不会替其他指针清除旧地址。
    \end{enumerate}
    \alert{局部指针离开作用域，不会自动delete它指向的动态对象。}
\end{frame}

\begin{frame}{malloc/free与new/delete差的不只是拼写}
    \renewcommand{\arraystretch}{1.5}
    \begin{tabular}{p{0.21\textwidth}>{\raggedright\arraybackslash}p{0.32\textwidth}>{\raggedright\arraybackslash}p{0.34\textwidth}}
        \textbf{对比项} & \textbf{malloc/free} & \textbf{new/delete表达式}\\
        \hline
        提供什么 & 原始存储，返回\mbox{\texttt{void *}} & 指定类型对象，返回\mbox{\texttt{T *}}\\
        创建类对象 & \textbf{不调用构造函数} & \texttt{new T(...)}调用构造函数\\
        结束类对象 & \textbf{不调用析构函数} & \texttt{delete}调用析构函数后释放\\
        普通分配失败 & \texttt{malloc}返回空指针 & 普通\texttt{new}抛出\texttt{std::bad\_alloc}\\
        配对要求 & \texttt{malloc}配\texttt{free} & \texttt{new}配\texttt{delete}
    \end{tabular}
    \par\medskip
    \alert{Student含有字符串，不能把malloc得到的地址直接当作构造好的Student使用。}\\
    \small 不混用分配与释放方式；不把\texttt{new}的结果交给\texttt{free}，反之亦然。
\end{frame}

\begin{frame}{数组用new[]创建，用delete[]释放}
    \begin{Code}*{cpp}
        \lstinputlisting{code/array.inc}
    \end{Code}
    \begin{itemize}
        \item \texttt{new int[2]\{90, 95\}}创建两个整数，初值分别为90和95。
        \item \texttt{scores}指向第一个元素；输出为\texttt{90 95}。
        \item \texttt{delete[] scores}释放整个数组；若元素为类对象，还会逐个调用析构函数。
    \end{itemize}
    \medskip
    \begin{vspquotered}{配对规则}
        单个对象：\texttt{new}配\texttt{delete}。\\[5pt]
        动态数组：\texttt{new[]}配\texttt{delete[]}，不能省去方括号。
    \end{vspquotered}
\end{frame}

    \VSPendframe{}
\end{document}
